% !PRINTED none
% section-opening page: "Siebenter Abschnitt. Realität der Wurzeln." heading block with §.76;
% no printed folio number was visible on this page in the scan (section-opening layout
% replaces the usual running head with the centred section title). Footer furniture
% "Weber, Algebra. I." and signature "16" are excluded under standing policy.

\begin{center}
\gesperrt{Seventh Section.}

\gesperrt{Reality of the Roots.}

\rule{3cm}{0.4pt}

\S.\,76.

\gesperrt{Generalities on the reality of roots of equations and on the discriminant.}
\end{center}

In this section we shall occupy ourselves with the question of how many roots of an algebraic equation are real. The coefficients of the equation are thereby assumed to be \gesperrt{real numbers}. We shall call such equations, for short, \gesperrt{real equations}, and begin with some general considerations.

The real equation of degree $n$
\[
f(x) = 0
\]
has, as we have seen in the third section, $n$ roots, which are either real or imaginary, i.\,e. of the form $\xi + i\eta$, with $\xi$, $\eta$ real. But this count of $n$ roots is obtained in general only if, in certain circumstances, we count a root multiply, namely $(m+1)$-fold when, together with $f(x)$, the first $m$ derivatives of $f(x)$ vanish. Now since a complex expression of the form $X + iY$ is equal to zero only when the two real constituents $X$ and $Y$ vanish separately, that is, when together with $X+iY$ also $X-iY$ vanishes, and since further, with real coefficients, if $f(\xi+i\eta) = X+iY$, one obtains $f(\xi-i\eta) = X-iY$, it follows from $f(\xi+i\eta)=0$ that at the same time $f(\xi-i\eta)=0$ must hold, and hence that to every imaginary root $\xi+i\eta$ there belongs a second imaginary root distinct from it, $\xi-i\eta$, i.\,e. the conjugate imaginary root
%CONT>>
